%
% Differentialkörper von Funktionen
%
\subsection{Differentialkörper von Funktionen}
In einem Funktionsausdruck für eine differenzierbare Funktion spielt
unter das Funktionsargument unter allen Symbolen 
eine besondere Rolle.
Wir gehen davon aus, dass man nach dem Funktionsargument ableiten
kann, weil es ``variabel'' ist.
Nach den Parametern gedenkt man dagegen nicht abzuleiten, da man
diese als ``konstant'' ansehen möchte.

%
% Derivation
%
\subsubsection{Derivation}
Im rein algebraischen Framework eines Funktionenkörpers ist es nicht
sinnvoll, die Definition der Ableitung als Grenzwert
\[
f'(z)
=
\lim_{h\to 0}
\frac{f(z+h)-f(z)}{h}
\]
eines Differenzenquotienten zu verwenden.
Um den Grenzwert zu definieren, wird ein Ordnungsrelation gebraucht,
ein Konzept, welches sich nicht mit rein algebraischen Eigenschaften
erfassen lässt.
Für den Grenzwert wird ein Zahlenkörper gebraucht, in dem alle
Grenzwerte existieren, was in $\mathbb{Q}$ nicht gegeben ist.
Der einzige Körper, der vollständig wäre, ist $\mathbb{R}$, der aber
im Computer nicht exakt dargestellt werden kann und daher nicht als
Basis eines CAS definiert werden kann.

Als alternativer Ansatz für die Definition der können die einfachsten
bekannten Rechenregeln der Ableitung verwendet werden.
Die Ableitung ist zum Beispiel eine lineare Abbildung.
Dies allein reicht jedoch nicht, weil die Ableitung auch für 
Potenzen $z^n$ oder für Produkte von Funktionen angewendet werden
soll.
Die Linearität kann darüber nichts aussagen.
Daher braucht es eine zusätzlich Regel für Produkte, die natürlich
die Produktregel für dei Ableitung wiedergeben soll.

\begin{definition}[Derivation]
Eine \emph{Derivation} in einem Funktionenkörper $K$ ist eine 
lineare Abbildung $D\colon K\to K$, die zusätzlich die Produktregel
\[
D(fg) = (Df)g + f(Dg)\qquad\forall f,g\in K
\]
erfüllt.
\index{Derivation}%
\end{definition}

In der Definition steht nichts davon, nach welcher Variablen da
abgeleitet werden soll.
In der Analysis verwendet man die Notation
\[
\frac{\partial}{\partial a},\quad
\frac{\partial}{\partial b},\quad
\frac{\partial}{\partial c},\quad
\frac{\partial}{\partial x},\quad
\frac{\partial}{\partial y}\qquad\text{oder}\qquad
\frac{d}{dx},\quad\text{und}\quad
\frac{d}{dz},
\]
die die Variable, nach der abgeleitet wird, klar erkennbar macht.
Allerdings ist in der Analysis auch die Notation $f'$ für die
abgeleitete Funktion üblich, aus der man die Variable auch nicht
ablesen kann.

Die spezielle Kennzeichnung des Funktionsargumentes ist aber auch
nicht nötig.
Betrachten wir $x$ als das Funktionsargument, dann ist $f(x)=x$ die
einfachste nichtkonstante Funktion, die wir uns denken können.
Die Ableitung dieser Funktion nach $x$ ist in der klassischen Analysis
ist $f'(x)=1$, wir können daher das Funktionsargument im Funktionenkörper
durch die Relation $Dx=1$ charakterisieren.

\begin{definition}[Differentialkörper]
\index{Differentialkorper@Differentialkörper}%
Ein \emph{Differentialkörper} $K$ ist ein Funktionenkörper mit einer
Derivation.
\end{definition}

Der Hauptsatz der Infinitesimalrechnung besagt, dass das Integral
einer Funktion auch eine Stammfunktion ist.
Die Definition des Integrals ist daher für die Arbeit eines
Computeralgebrasystems nicht relevant, es mur nur eine Stammfunktion
finden können.
Der klassische Begriff der Stammfunktion verwendet die Ableitung.
Für den vorliegenden algebraischen Zusammenhang kann er wie folgt
verallgemeinert werden.

\begin{definition}[Stammfunktion]
Sei $F$ ein Differentialkörper mit Derivation $D$.
Eine Funktion $\vartheta\in F$ heisst eine \emph{Stammfunktion} von 
$f\in F$, wenn $D\vartheta=f$.
\index{Stammfunktion}%
\end{definition}

%
% Ableitungsregeln
%
\subsubsection{Ableitungsregeln}
Die Produktregel reicht bereits, um die aus dem Analysisunterricht
bekannten Ableitungsregeln zu reproduzieren.

\begin{enumerate}
\item Ableitung einer Potenz: wir betrachten die Potenz $f^n$ eines
Elementes $f\in K$.
Die Derivationseigenschaft liefert
\begin{align*}
D(f^n)
&=
D(\underbrace{f\cdot f\cdot\ldots\cdot f}_{\displaystyle n})
\\
&=
Df\cdot\ldots f
+
f\cdot (Df)\cdot\ldots \cdot f
+
\ldots
+
f\cdot f \cdot\ldots \cdot Df
=
nf^{n-1}\cdot Df.
\end{align*}
Die ist die aus der Analysis wohlbekannte Regel für Ableitungen einer
Potenz.
\item Ableitung der Eins: Die Eins zeichnet sich durch die
Eigenschaft $1^n=1$ aus.
Lässt man die Derivation darauf wirken, erhält man
\[
D(1^n)
=
n1^{n-1}D1
=
D1
\qquad\Rightarrow\qquad
(n-1)D1 = 0.
\]
Da $n-1\ne 1$ ist, folgt $D1=0$, die Ableitung der Eins ist somit $1$.
Da die Derivation $D$ linear ist, gilt für jedes andere Element
$q\in\mathbb{Q}$ sofort $Dq = D(q\cdot 1) = q\cdot D1 = 0$.
\item
Ableitung von $1/f$: Die reziproke Funktion $1/f$ ist charakterisiert
durch die Eigenschaft $f\cdot (1/f)=1$.
Die Ableitung davon ist
\begin{align*}
&&
D(f\cdot(1/f))
&=
(Df)\cdot (1/f) + f\cdot D(1/f)
=
D1
=
0
\\
&\Rightarrow&
f\cdot D(1/f) &= -\frac{Df}{f}
\\
&\Rightarrow&
D(1/f)
&=
-\frac{Df}{f^2},
\end{align*}
dies ist die Ableitungsregel einer Potenz $f^{-1}$.
\item Potenzregel für beliebige Exponenten: Wir wissen bereits,
dass $D(f^n)=nf^{n-1}Df$ für positive $n$.
Für einen negativen Exponenten ist
\[
D(f^{-n})
=
D\biggl(\frac{1}{f^n}\biggr)
=
-
\frac{D(f^n)}{f^{2n}}
=
-\frac{nf^{n-1}Df}{f^{2n}}
=
-
n
f^{-n-1}
Df,
\]
dies ist die erwartete Ableitungsformel für beliebige ganzzahlige 
Exponenten.
\item Ableitungsregel für Polynome:
Sie $p(X) = p(X) = p_0 + p_1X + \dots + p_nX^n \in K[X]$ ein Polynom
mit Koeffizienten im Koeffizientenkörper.
Wir möchten die Ableitung von $Dp(f)$ berechnen.
Nach der Potenzregel gilt
\begin{align*}
Dp(f)
&=
D(p_0+p_1f + p_2f^2 + \ldots + p_{n-1}f^{n-1}+p_nf^n)
\\
&=
p_1Df + p_2\cdot 2f Df + \ldots p_{n-1}(n-1)f^{n-2}Df + p_nnf^{n-1}Df
\\
&=
\bigl(
p_1 + 2p_2 f + 3p_3f^2 + \ldots + p_{n-1}(n-1)f^{n-2} + p_nnf^{n-1}
\bigr)
Df.
\intertext{Das Polynom auf der rechten Seite ist auch bekannt als
das abgeleitete Polynom $p'(X)$, so dass man}
&= p'(f)\, Df
\end{align*}
schreiben kann.
Dies ist formal die Kettenregel für Polynome.
\item Quotientenregel:
Der Quotient $f/g$ kann auch als Produkt $f\cdot (1/g)$ geschrieben
werden.
Daher ist die Ableitung davon
\begin{align*}
D\frac{f}{g}
&=
D\biggl(f\cdot\frac1g\biggr)
=
Df\cdot\frac1g + f\cdot D\frac1g
\\
&=
Df\cdot\frac1g + f\cdot \biggl(-\frac{Dg}{g^2}\biggr)
\\
&=
\frac{Df}{g} - \frac{f\cdot Dg}{g^2}
=
\frac{Df\cdot g - f\cdot Dg}{g^2}.
\end{align*}
Dies ist die bekannte Quotientenregel.
\item Ableitung einer rationalen Funktion:
Eine rationale Funktion ist ein Quotient $p[X]/q[X]$ mit
Polynomen $p,q\in K[X]$.
Wir möchten $p(f)/q(f)$ ableiten:
\begin{align*}
D\biggl(\frac{p(f)}{q(f)}\biggr)
&=
\frac{
(Dp(f))\cdot q(f) - p(f)\cdot Dq(f)
}{
q(f)^2
}
\\
&=
\frac{p'(f)\,Df\cdot q(f)-p(f)\cdot q'(f)\, Df}{q(f)^2}
\\
&=
\frac{p'(f)q(f)-p(f) q'(f)}{q(f)^2}
Df.
\end{align*}
Dies ist die Kettenregel für den Polynombruch $p(X)/q(X)$.
\end{enumerate}

Diese Beispiele zeigen, dass alle Regeln, die man üblicherweise
zur Berechnung von Ableitungen braucht, allein aus der Linearität
und Eigenschaft der Derivation folgen.
Der Funktionenkörper $F$ mit der Derivation $D$ ist daher ein
vollständiges Modell für die algebraischen Eigenschaften der
Ableitung.
Es ist daher auch denkbar, dass das Problem, eine Stammfunktion
zu finden, allein in diesem algebraischen Rahmen gelöst werden
kann.

Der Vollständigkeit halber halten wir hier noch die Definition
des abgeleiteten Polynoms fest.

\begin{definition}[Ableitung eines Polynoms]
\index{Ableitung eines Polynoms}%
Ist $p(X)\in K[X]$ ein Polynom mit Koeffizienten in einem
Koeffizientenkörper mit der Darstellung
\[
p(X)
=
p_0 + p_1X + p_2X^2 + \ldots + p_{n-1}X^{n-1} + p_nX^n,
=
\sum_{k=0}^n p_k X^k
\]
dann heisst
\[
p'(X)
=
p_1 + 2p_2 X + 3p_3 X^2 + \ldots + (n-1)p_{n-1}X^{n-2} + np_n X^{n-1}
=
\sum_{k=1}^n kp_k X^{k-1}
\]
das \emph{abgeleitete Polynom}.
\end{definition}

%
% Ableitung eines algebraischen Elementes
%
\subsubsection{Ableitung eines algebraischen Elementes}
Wenn der Funktionenkörper $F$ um ein algebraisches Element $\alpha$
erweitert wird, dann muss auch die Ableitung $D\alpha$ berechnet
werden können, damit $F(\alpha)$ wieder ein Differentialkörper ist.

Sei als $p(X)\in F[X]$ das Minimalpolynom des Elementes $\alpha$,
es gilt also $p(\alpha)=0$.
Die Ableitung ist nach den Berechungsregeln
\begin{align*}
0
=
Dp(\alpha)
&=
p'(\alpha)D\alpha
+
(Dp_0+Dp_1\alpha + \dots + Dp_n\alpha^n)
\intertext{Da $p$ das Minimalpolynom von $\alpha$ ist, kann $p'(\alpha)$
nicht verschwinden, denn dann wäre $p'$ ein Polynom kleineren Grades,
das auf $\alpha$ verschwindet.
Somit kann man die Gleichung nach $D\alpha$ auflösen und erhält}
D\alpha
&=
-
\frac{Dp_0+(Dp_1)\alpha +\ldots+(Dp_n)\alpha^n}{p'(\alpha)}.
\end{align*}
Die Ableitung eines algebraischen Elementes ist also durch das
Minimalpolynom eindeutig bestimmt.

\begin{beispiel}
Die Quadratwurzel $\!\sqrt{f}$ eines Elementes $f\in F$ hat das
Minimalpolynom $p(X) = X^2 -f \in F[X]$ mit der Ableitung $p'(X)=2X$.
Die Ableitung der Koeffizienten ergibt das Polynom $-Df$.
Damit ist die Ableitung der Quadratwurzel
\begin{align*}
D\!\sqrt{f}
&=
\frac{-Df}{p'(\!\sqrt{f})}
=
-\frac{1}{2\!\sqrt{f}}\cdot Df.
\end{align*}
Dies ist die Kettenregel für die Wurzelfunktion.
\end{beispiel}

\begin{beispiel}
Die $n$-te Wurzel $\sqrt[n]{f}$ eines Elementes $f\in K$ ist algebraisch
mit dem Minimalpolynom $p(X)=X^n -f$ mit Ableitung $p'(X)=nX^{n-1}$ und
dem Polynom mit den abgeleiteten Koeffizienten $-Df$.
Dann ist
\begin{align*}
D\!\sqrt[n]{f}
&=
\frac{-Df}{p'\bigl(\!\sqrt[n]{f}\bigr)}
=
-
\frac{1}{n\bigl(\!\sqrt[n]{f}\bigr)^{n-1}}
\cdot
Df.
\end{align*}
Dies ist die Kettenregel für die $n$-Wurzel.
\end{beispiel}

%
% Konstantenkörper
%
\subsubsection{Konstantenkörper}
Bei der Konstruktion des Funktionskörpers sind wir davon ausgegangen,
dass wir bereits wissen, welche Objekte Konstanten sind.
Die Ableitung $D$ ermöglicht jetzt aber auch, die Konstanten dadurch
zu definieren, dass ihre Ableitung verschwindet.

\begin{definition}[Konstantenkörper]
Sei $F$ ein Differentialkörper, dann heisst
\[
C=\{c\in F\mid Dc=0\}
\]
der Körper der Konstanten oder \emph{Konstantenkörper} in $F$.
\index{Konstantenkorper@Konstantenkörper}%
\end{definition}

Wir müssen uns noch überlegen, dass $C$ wirklich ein Körper ist.
Dazu müssen wir nachrechnen, dass Summe, Produkt und Quotient
konstanter Elemente aus $C$ wieder in $C$ liegt.
Wir wählen daher $a,b\in C$, d.~h.~$Da=0$ und $Db=0$.
Dann gilt auch
\begin{align*}
D(a\pm b)
&=
Da\pm Db = 0\pm 0 = 0
&&\Rightarrow&
a\pm b&\in C
\\
D(ab)
&=
(Da)\cdot b + a\cdot (Db)
=
0\cdot b + a\cdot 0
=
0
&&\Rightarrow&
ab&\in C
\\
D\frac{a}{b}
&=
\frac{(Da)\cdot b - a\cdot (Db)}{b^2}
=
\frac{0\cdot b - a\cdot 0}{b^2}
=
0.
&&\Rightarrow&
\frac{a}{b}&\in C.
\end{align*}
$C$ ist somit ein Körper.
Der Konstantenkörper kann grösser sein als der ursprüngliche
Koeffizientenkörper.
